\section{Symmetry of Quadrupolar Interactions between Tb$^{3+}$ ions in TbPO$_4$}

The symmetry of quadrupolar interactions between Tb$^{3+}$ ions in TbPO$_4$ is investigated
up to the third nearest neighbour using group theory.
 For the first and second nearest neighbour
13 independent terms are allowed, while for the third nearest neighbour  25 independent terms
are possible. 

In order to calculate the dispersion of excitations in $TbPO_4$ a careful analysis
of the symmetry of the quadrupolar interactions is necessary. In this report I attempt
to apply group theory in order to derive the allowed quadrupolar terms for interactions
up to the third nearest neighbour.
The general form of a quadrupolar interaction 
in the Hamiltonian is 

\begin{equation}
H_Q=-\frac{1}{2}\sum_{abmm'}K(ambm') O_{2m}(\mbf J_a)O_{2m'}(\mbf J_b)
\end{equation}

Here
$O_{lm}(\mbf J_a)$ corresponds to the Stevens Operator~\cite{hutchings64-227}
for the Tb atom $a$ with total angular momentum $\mbf J_a$. Quadrupolar terms with
$m=m'$ will contribute to the $K^{\mu}$
 in equation (4) in Forino et al.2026:

\begin{equation}
K^{\mu}=\sum_b K(ambm) \qquad (\mu=\alpha,\gamma,\delta)
\end{equation}

where $m=0,2,-2$ for $\mu=\alpha,\gamma,\delta$, respectively and

\begin{equation}
K^{\epsilon}=\sum_{b,m=1,-1} K(ambm)
\end{equation}

We will show, that also $K(ambm')$ with $m\neq m'$ can be nonzero. However
these will not contribute to $K^{\mu}$, because contributions from 
the various neighbours cancel each other. These interactions
 with $m\neq m'$ cannot be seen in bulk properties of $TbPO_4$, because
$TbPO_4$ exhibits ferroquadrupolar order. Therefore I call these
interactions ''hidden'' interactions. Hidden quadrupolar interactions
may influence the dispersion of excitations. For sake of simplicity,
we have (up to now) not considered these interactions.

The types of Tb-neighbours of Tb at $(0,3/4,1/8)$ are given in table~\ref{nn}. In the 
following sections the symmetry relations of quadrupolar interactions
will be derived based on the numbering in this table. 

In section \ref{13ind} the 13 independent terms in the nearest and second
nearest neighbour interactions are listed, in section \ref{25ind} the
25 independent terms in the third nearst neighbour interactions are given.
The interactions with $m=m'=0,2,-2$ are boring, these are equal for equivalent 
neighbours. However, looking at $m=m'=\pm1$ yields an important
 result: the relations of $K(a1b1)$ and
$K(a,-1,b,-1)$ for the various neighbours,  see equations
(\ref{nn211}),(\ref{nn121}),(\ref{nnn211}),(\ref{nnn121}).


\begin{table}
\caption{Types of Tb-neighbours in $TbPO_4$, when neighbours for the 
two Tb  sites differ, the upper sign refers to neighbours of Tb at
$(0,3/4,1/8)$  and the lower sign refers to neighbours of Tb at $(1/2,3/4,3/8)$}
\label{nn}\begin{tabular}{c|cc}
type & b & relative coordinates \\
\hline
nearest  & 1 & $(+1/2,0,\pm1/4)$ \\
neighbour & 2 & $(-1/2,0,\pm1/4)$ \\
          & 3 & $(0,+1/2,\mp1/4)$ \\
          & 4 & $(0,-1/2,\mp1/4)$ \\
\hline
next   & 1 & $(+1/2,0,\mp3/4)$ \\
nearest & 2 & $(-1/2,0,\mp3/4)$ \\
neighbour & 3 & $(0,+1/2,\pm3/4)$ \\
          & 4 & $(0,-1/2,\pm3/4)$ \\
\hline
third     & 1  & $(+1/2,+1/2,+1/2)$ \\
nearest   & 2  & $(+1/2,+1/2,-1/2)$ \\
neighbour & 3  & $(+1/2,-1/2,+1/2)$ \\
          & 4  & $(+1/2,-1/2,-1/2)$ \\
          & 5  & $(-1/2,+1/2,+1/2)$ \\
          & 6  & $(-1/2,+1/2,-1/2)$ \\
          & 7  & $(-1/2,-1/2,+1/2)$ \\
          & 8  & $(-1/2,-1/2,-1/2)$ \\
\end{tabular}
\end{table}

\subsection{Point Group of Tb in TbPO$_4$}

The tetragonal space group of TbPO$_4$ is $I 4_1/amd$ ($D_{4h}^{19}$, number 141 in international 
tables). With origin choice 2 the Tb occupies the $4a$ sites at $(0,3/4,1/8)$ and
$(1/2,3/4,3/8)$. The point group symmetry on the Tb site is 
 $\bar 4 m 2 $ ($D_{2d}$). This point group has 8 elements. The character table
of the irreducible representations is given in table~\ref{char}.

\begin{table}
\caption{Character table of irreducible representations of point group $\bar 4 m 2 $ ($D_{2d}$).}
\label{char}\begin{tabular}{c|ccccc}
$D_{2d}$ & $E$ & $C_2$ & $2S_4$ & $ 2C'_2$ & $2\sigma_d$ \\
\hline
$A_1$    &  1  & 1     & 1      & 1        &   1   \\
$A_2$    &  1  & 1     & 1      & -1        &   -1   \\
$B_1$    &  1  & 1     & -1      & 1        &   1   \\
$B_2$    &  1  & 1     & -1      & -1        &   1   \\
$E$      &  2  & -2    & 0       & 0         &   0 \\
\end{tabular}
\end{table}





% ------------------------------------------------------------------------------


\subsection{The reducible representation $T^1(G)$ for nearest  and next nearest Tb-Tb neighbours}

The nearest and next nearest neighbour interactions transform according to the
same reducible representation $T^1(G)$. We will construct this reducible representation 
and use it's characters $\chi_p^{T^1}$ to determine how often the irreducible unit representation
$A_1$ occurs in it's reduction. Moreover we will construct the projection operator
$P^{A_1}$ to find the basis vectors of the corresponding irreducible subspaces. 
These basis vectors correspond to the independent terms in the quadrupolar Hamiltonian.

Without considering symmetry a general term in the quadrupolar Tb-Tb interaction
will have the form $O_{2m}(\mbf J)O_{2m'}(\mbf J_b))$. Here $\mbf J$ denotes the
total angular momentum at the first Tb site, and $\mbf J_b$ the total angular
momentum on it's neighbour number $b$ according to the table~\ref{nn}. 
In order to apply group theory to find symmetry allowed interactions we
can think of this general term as on of the basis vectors
 for a representation $T^1(G)$ of the point groupo $\bar 4 m 2 $. We will
use the notation $|mm'b\rangle$ to denote $O_{2m}(\mbf J)O_{2m'}(\mbf J_b)$.
We will in the following sections determine the representation matrices 
$T^1_{m_1m_2b,m'_1m'_2b'}(G)$ which describe the transformation
of the state $|m_1m_2b\rangle$ when applying a point symmetry 
operation $G$:

\begin{equation}
|m_1m_2b\rangle \stackrel{G}{\rightarrow} \sum_{m'_1m'_2b'} T^1_{m_1m_2b,m'_1m'_2b'}(G)|m'_1m'_2b'\rangle
\end{equation}
% ------------------------------------------------------------------------------

\subsubsection{The group element $E$}
 
For the identity $E$ the representation 

\begin{equation}
T^1_{m_1m_2b,m'_1m'_2b'}(E)=\delta_{m_1,m'_1}\delta_{m_2,m'_2}\delta_{bb'}
\end{equation}

Writing the $bb'$-dependence in Matrix notation we get

\begin{equation}\label{E}
T^1_{m_1m_2,m'_1m'_2}(E)=\delta_{m_1,m'_1}\delta_{m_2,m'_2}
\left (
\begin{array}{cccc}
1 & 0 & 0 & 0 \\
0 & 1 & 0 & 0 \\
0 & 0 & 1 & 0 \\
0 & 0 & 0 & 1 
\end{array} \right )
\end{equation}


The character is given by

\begin{equation}
\chi^{T^1}(E)=\Trace{T^1(E)}=5\times5\times4=100
\end{equation}

% ------------------------------------------------------------------------------

\subsubsection{The group element $C_2$}

$C_2$ is a two fold axis along $z$, this corresponds to 

\begin{eqnarray}
(x,y,z) &\stackrel{C_2}{\rightarrow}& (-x,-y,z) \nn
(J_x,J_y,J_z) &\stackrel{C_2}{\rightarrow}&  (-J_x,-J_y,J_z) \nn
J_\pm &\stackrel{C_2}{\rightarrow}& - J_\pm \nn
\end{eqnarray}

We will use the definition of Stevens Operators:

\begin{eqnarray}
O^s_{22}=O_{2-2}&=&\frac{-i}{2}[J_+^2-J_-^2]=J_xJ´_y+J_yJ_x=2P_{xy}=I_4 \nn
O^s_{21}=O_{2-1}&=&\frac{-i}{4}[J_z(J_+-J_-)+(J_+-J_-)J_z]=\frac{1}{2}[J_yJ_z+J_zJ_y]=P_{yz}=I_5 \nn
O_{20}&=&[3J_z^2-X]=I_6 \nn
O_{21}&=&\frac{1}{4}[J_z(J_++J_-)+(J_++J_-)J_z]=\frac{1}{2}[J_xJ_z+J_zJ_x]=P_{xz}=I_7 \nn
O_{22}&=&\frac{1}{2}[J_+^2+J_-^2]=J_x^2-J_y^2 =I_8\nn
\end{eqnarray}

Thus 
\begin{eqnarray}
O_{2-2}&\stackrel{C_2}{\rightarrow}&O_{2-2} \nn
O_{2-1}&\stackrel{C_2}{\rightarrow}&-O_{2-1} \nn
O_{20}&\stackrel{C_2}{\rightarrow}&O_{20} \nn
O_{21}&\stackrel{C_2}{\rightarrow}&-O_{21}\nn
O_{22}&\stackrel{C_2}{\rightarrow}&O_{22}\nn
\end{eqnarray}

or in summary
\begin{equation}
O_{2m}\stackrel{C_2}{\rightarrow} (-1)^m O_{2m} 
\end{equation}

and 

\begin{equation}
T^1_{m_1m_2b,m'_1m'_2b'}(C_2)=(-1)^{m_1+m_2}\delta_{m_1,m'_1}\delta_{m_2,m'_2}
(\delta_{b1}\delta_{b'2}+\delta_{b2}\delta_{b'1}+\delta_{b3}\delta_{b'4}+\delta_{b4}\delta_{b'3})
\end{equation}

Writing the $bb'$-dependence in Matrix notation we get

\begin{equation}\label{C2}
T^1_{m_1m_2,m'_1m'_2}(C_2)=(-1)^{m_1+m_2}\delta_{m_1,m'_1}\delta_{m_2,m'_2}
\left (
\begin{array}{cccc}
0 & 1 & 0 & 0 \\
1 & 0 & 0 & 0 \\
0 & 0 & 0 & 1 \\
0 & 0 & 1 & 0 
\end{array} \right )
\end{equation}

The character is given by the trace and is zero

\begin{equation}
\chi^{T^1}(C_2)=\Trace{T^1(C_2)}=0
\end{equation}
% ------------------------------------------------------------------------------

\subsubsection{The clockwise mirror rotation $S_4$}

 The clockwise mirror rotation $S_4$ can be written as $S_4=\sigma C_4=IC_2C_4$ ($\sigma$ is 
a horizontal mirror plane).

\begin{eqnarray}
(x,y,z) &\stackrel{S_4}{\rightarrow}& (y,-x,-z) \nn
(J_x,J_y,J_z) &\stackrel{S_4}{\rightarrow}&  (-J_y,J_x,J_z) \nn
J_\pm &\stackrel{S_4}{\rightarrow}& \pm iJ_\pm \nn
\end{eqnarray}


\begin{eqnarray}
´O_{2-2}&\stackrel{S_4}{\rightarrow}&-O_{2-2} \nn
O_{2-1}&\stackrel{S_4}{\rightarrow}&O_{21} \nn
O_{20}&\stackrel{S_4}{\rightarrow}&O_{20} \nn
O_{21}&\stackrel{S_4}{\rightarrow}&-O_{2-1}\nn
O_{22}&\stackrel{S_4}{\rightarrow}&-O_{22}\nn
\end{eqnarray}

or in summary
\begin{equation}
O_{2m}\stackrel{S_4}{\rightarrow} 
(\delta_{m0}-\delta_{m1}+\delta_{m,-1}-\delta_{m2}-\delta_{m,-2})
 O_{2,(-1)^mm} 
\end{equation}

and 

\begin{eqnarray}
T^1_{m_1m_2b,m'_1m'_2b'}(S_4)&=&
(\delta_{m_10}+\delta_{m_11}-\delta_{m_1,-1}-\delta_{m_12}-\delta_{m_1,-2})
\delta_{(-1)^{m_1}m_1,m'_1} \nn
&&
(\delta_{m_20}+\delta_{m_21}-\delta_{m_2,-1}-\delta_{m_22}-\delta_{m_2,-2})
\delta_{(-1)^{m_2}m_2,m'_2} \nn
&&(\delta_{b1}\delta_{b'4}+\delta_{b2}\delta_{b'3}+\delta_{b3}\delta_{b'1}+\delta_{b4}\delta_{b'2})
\end{eqnarray}

Writing the $bb'$-dependence in Matrix notation we get

\begin{eqnarray}\label{S4}
T^1_{m_1m_2,m'_1m'_2}(S_4)&=&(\delta_{m_10}+\delta_{m_11}-\delta_{m_1,-1}-\delta_{m_12}-\delta_{m_1,-2})
\delta_{(-1)^{m_1}m_1,m'_1} \nn
&&
(\delta_{m_20}+\delta_{m_21}-\delta_{m_2,-1}-\delta_{m_22}-\delta_{m_2,-2})
\delta_{(-1)^{m_2}m_2,m'_2} \nn
&&\left (
\begin{array}{cccc}
0 & 0 & 0 & 1 \\
0 & 0 & 1 & 0 \\
1 & 0 & 0 & 0 \\
0 & 1 & 0 & 0 
\end{array} \right )
\end{eqnarray}

The character is given by the trace and is zero

\begin{equation}
\chi^{T^1}(S_4)=\Trace{T^1(S_4)}=0
\end{equation}

% ------------------------------------------------------------------------------

\subsubsection{The counter-clockwise mirror rotation $\bar S_4$}

 The counterclockwise mirror rotation $\bar S_4$ can be written as $\bar S_4=\sigma \bar C_4=I\bar C_2 \bar C_4$ ($\sigma$ is 
a horizontal mirror plane).

\begin{eqnarray}
(x,y,z) &\stackrel{\bar S_4}{\rightarrow}& (-y,x,-z) \nn
(J_x,J_y,J_z) &\stackrel{\bar S_4}{\rightarrow}&  (J_y,-J_x,J_z) \nn
J_\pm &\stackrel{\bar S_4}{\rightarrow}& \mp iJ_\pm \nn
\end{eqnarray}

\begin{eqnarray}
´O_{2-2}&\stackrel{\bar S_4}{\rightarrow}&-O_{2-2} \nn
O_{2-1}&\stackrel{\bar S_4}{\rightarrow}&-O_{21} \nn
O_{20}&\stackrel{\bar S_4}{\rightarrow}&O_{20} \nn
O_{21}&\stackrel{\bar S_4}{\rightarrow}&O_{2-1}\nn
O_{22}&\stackrel{\bar S_4}{\rightarrow}&-O_{22}\nn
\end{eqnarray}

or in summary
\begin{equation}
O_{2m}\stackrel{\bar S_4}{\rightarrow} 
\{(i)^m [1-(m \bmod 2)] +m(m \bmod 2) \} 
(\delta_{m0}-\delta_{m1}+\delta_{m,-1}-\delta_{m2}-\delta_{m,-2})
O_{2,(-1)^mm} 
\end{equation}

and 

\begin{eqnarray}
T^1_{m_1m_2b,m'_1m'_2b'}(\bar S_4)&=&
(\delta_{m_10}-\delta_{m_11}+\delta_{m_1,-1}-\delta_{m_12}-\delta_{m_1,-2})
\delta_{(-1)^{m_1}m_1,m'_1} \nn
&&
(\delta_{m_20}-\delta_{m_21}+\delta_{m_2,-1}-\delta_{m_22}-\delta_{m_2,-2})
\delta_{(-1)^{m_2}m_2,m'_2} \nn
&&(\delta_{b1}\delta_{b'3}+\delta_{b2}\delta_{b'4}+\delta_{b3}\delta_{b'2}+\delta_{b4}\delta_{b'1})
\end{eqnarray}

Writing the $bb'$-dependence in Matrix notation we get

\begin{eqnarray}\label{bS4}
T^1_{m_1m_2,m'_1m'_2}(\bar S_4)&=&
(\delta_{m_10}-\delta_{m_11}+\delta_{m_1,-1}-\delta_{m_12}-\delta_{m_1,-2})
\delta_{(-1)^{m_1}m_1,m'_1} \nn
&&
(\delta_{m_20}-\delta_{m_21}+\delta_{m_2,-1}-\delta_{m_22}-\delta_{m_2,-2})
\delta_{(-1)^{m_2}m_2,m'_2} \nn
&&\left (
\begin{array}{cccc}
0 & 0 & 1 & 0 \\
0 & 0 & 0 & 1 \\
0 & 1 & 0 & 0 \\
1 & 0 & 0 & 0 
\end{array} \right )
\end{eqnarray}

The character is given by the trace and is zero

\begin{equation}
\chi^{T^1}(\bar S_4)=\Trace{T^1(\bar S_4)}=0
\end{equation}

% ------------------------------------------------------------------------------

\subsubsection{The two fold horizontal [110]-axis  $C'_2$}

 The    [110]-axis is a two fold axis $C'_2$ and the transformation is given by

\begin{eqnarray}
(x,y,z) &\stackrel{C'_2}{\rightarrow}& (y,x,-z) \nn
(J_x,J_y,J_z) &\stackrel{C'_2}{\rightarrow}&  (J_y,J_x,-J_z) \nn
J_\pm &\stackrel{C'_2}{\rightarrow}& \pm i J_\mp \nn
\end{eqnarray}

\begin{eqnarray}
´O_{2-2}&\stackrel{C'_2}{\rightarrow}&O_{2-2} \nn
O_{2-1}&\stackrel{C'_2}{\rightarrow}&-O_{21} \nn
O_{20}&\stackrel{C'_2}{\rightarrow}&O_{20} \nn
O_{21}&\stackrel{C'_2}{\rightarrow}&-O_{2-1}\nn
O_{22}&\stackrel{C'_2}{\rightarrow}&-O_{22}\nn
\end{eqnarray}

or in summary
\begin{equation}
O_{2m}\stackrel{C'_2}{\rightarrow} (\delta_{m0}-\delta_{m1}-\delta_{m,-1}-\delta_{m2}+\delta_{m,-2}) O_{2,(-1)^mm} 
\end{equation}

and 

\begin{eqnarray}
T^1_{m_1m_2b,m'_1m'_2b'}(C'_2)&=&(\delta_{m_10}-\delta_{m_11}-\delta_{m_1,-1}-\delta_{m_12}+\delta_{m_1,-2})
\delta_{(-1)^{m_1}m_1,m'_1}\nn
&&(\delta_{m_20}-\delta_{m_21}-\delta_{m_2,-1}-\delta_{m_22}+\delta_{m_2,-2})
\delta_{(-1)^{m_2}m_2,m'_2}\nn
&&(\delta_{b1}\delta_{b'3}+\delta_{b2}\delta_{b'4}+\delta_{b3}\delta_{b'1}+\delta_{b4}\delta_{b'2})
\end{eqnarray}

Writing the $bb'$-dependence in Matrix notation we get

\begin{eqnarray}\label{Cd2}
T^1_{m_1m_2,m'_1m'_2}(C'_2)&=&(\delta_{m_10}-\delta_{m_11}-\delta_{m_1,-1}-\delta_{m_12}+\delta_{m_1,-2})
\delta_{(-1)^{m_1}m_1,m'_1}\nn
&&(\delta_{m_20}-\delta_{m_21}-\delta_{m_2,-1}-\delta_{m_22}+\delta_{m_2,-2})
\delta_{(-1)^{m_2}m_2,m'_2}\nn
&&\left (
\begin{array}{cccc}
0 & 0 & 1 & 0 \\
0 & 0 & 0 & 1 \\
1 & 0 & 0 & 0 \\
0 & 1 & 0 & 0 
\end{array} \right )
\end{eqnarray}

The character is given by the trace and is zero

\begin{equation}
\chi^{T^1}(C'_2)=\Trace{T^1(C'_2)}=0
\end{equation}

% ------------------------------------------------------------------------------

\subsubsection{The two fold horizontal $[1\bar 10]$-axis  $C''_2$}

 The    $[1\bar 10]$-axis is a two fold axis $C''_2$ and the transformation is given by

\begin{eqnarray}
(x,y,z) &\stackrel{C''_2}{\rightarrow}& (-y,-x,-z) \nn
(J_x,J_y,J_z) &\stackrel{C''_2}{\rightarrow}&  (-J_y,-J_x,-J_z) \nn
J_\pm &\stackrel{C''_2}{\rightarrow}& \mp i J_\mp \nn
\end{eqnarray}

\begin{eqnarray}
´O_{2-2}&\stackrel{C''_2}{\rightarrow}&O_{2-2} \nn
O_{2-1}&\stackrel{C''_2}{\rightarrow}&O_{21} \nn
O_{20}&\stackrel{C''_2}{\rightarrow}&O_{20} \nn
O_{21}&\stackrel{C''_2}{\rightarrow}&O_{2-1}\nn
O_{22}&\stackrel{C''_2}{\rightarrow}&-O_{22}\nn
\end{eqnarray}

or in summary
\begin{equation}
O_{2m}\stackrel{C''_2}{\rightarrow} (\delta_{m0}+\delta_{m1}+\delta_{m,-1}-\delta_{m2}+\delta_{m,-2}) O_{2,(-1)^mm} 
\end{equation}

and 

\begin{eqnarray}
T^1_{m_1m_2b,m'_1m'_2b'}(C''_2)&=&(\delta_{m_10}+\delta_{m_11}+\delta_{m_1,-1}-\delta_{m_12}+\delta_{m_1,-2})
\delta_{(-1)^{m_1}m_1,m'_1}\nn
&&(\delta_{m_20}+\delta_{m_21}+\delta_{m_2,-1}-\delta_{m_22}+\delta_{m_2,-2})
\delta_{(-1)^{m_2}m_2,m'_2}\nn
&&(\delta_{b1}\delta_{b'4}+\delta_{b2}\delta_{b'3}+\delta_{b3}\delta_{b'2}+\delta_{b4}\delta_{b'1})
\end{eqnarray}

Writing the $bb'$-dependence in Matrix notation we get

\begin{eqnarray}\label{Cdd2}
T^1_{m_1m_2,m'_1m'_2}(C''_2)&=&(\delta_{m_10}+\delta_{m_11}+\delta_{m_1,-1}-\delta_{m_12}+\delta_{m_1,-2})
\delta_{(-1)^{m_1}m_1,m'_1}\nn
&&(\delta_{m_20}+\delta_{m_21}+\delta_{m_2,-1}-\delta_{m_22}+\delta_{m_2,-2})
\delta_{(-1)^{m_2}m_2,m'_2}\nn
&&\left (
\begin{array}{cccc}
0 & 0 & 0 & 1 \\
0 & 0 & 1 & 0 \\
0 & 1 & 0 & 0 \\
1 & 0 & 0 & 0 
\end{array} \right )
\end{eqnarray}

The character is given by the trace and is zero

\begin{equation}
\chi^{T^1}(C''_2)=\Trace{T^1(C''_2)}=0
\end{equation}

% ------------------------------------------------------------------------------

\subsubsection{The yz mirror plane  $\sigma_d$}

For the yz mirror plane $\sigma_d$ the transformation is given by

\begin{eqnarray}
(x,y,z) &\stackrel{\sigma_d}{\rightarrow}& (-x,y,z) \nn
(J_x,J_y,J_z) &\stackrel{\sigma_d}{\rightarrow}&  (J_x,-J_y,-J_z) \nn
J_\pm &\stackrel{\sigma_d}{\rightarrow}&  J_\mp \nn
\end{eqnarray}

\begin{eqnarray}
´O_{2-2}&\stackrel{\sigma_d}{\rightarrow}&-O_{2-2} \nn
O_{2-1}&\stackrel{\sigma_d}{\rightarrow}&O_{2-1} \nn
O_{20}&\stackrel{\sigma_d}{\rightarrow}&O_{20} \nn
O_{21}&\stackrel{\sigma_d}{\rightarrow}&-O_{21}\nn
O_{22}&\stackrel{\sigma_d}{\rightarrow}&O_{22}\nn
\end{eqnarray}

or in summary
\begin{equation}
O_{2m}\stackrel{\sigma_d}{\rightarrow} (\delta_{m0}-\delta_{m1}+\delta_{m,-1}+\delta_{m2}-\delta_{m,-2}) O_{2,m} 
\end{equation}

and 

\begin{eqnarray}
T^1_{m_1m_2b,m'_1m'_2b'}(\sigma_d)&=&(\delta_{m_10}-\delta_{m_11}+\delta_{m_1,-1}+\delta_{m_12}-\delta_{m_1,-2})
\delta_{m_1,m'_1}\nn
&&(\delta_{m_20}-\delta_{m_21}+\delta_{m_2,-1}+\delta_{m_22}-\delta_{m_2,-2})
\delta_{m_2,m'_2}\nn
&&(\delta_{b1}\delta_{b'2}+\delta_{b2}\delta_{b'1}+\delta_{b3}\delta_{b'3}+\delta_{b4}\delta_{b'4})
\end{eqnarray}

Writing the $bb'$-dependence in Matrix notation we get

\begin{eqnarray}\label{sd}
T^1_{m_1m_2,m'_1m'_2}(\sigma_d)&=&(\delta_{m_10}-\delta_{m_11}+\delta_{m_1,-1}+\delta_{m_12}-\delta_{m_1,-2})
\delta_{m_1,m'_1}\nn
&&(\delta_{m_20}-\delta_{m_21}+\delta_{m_2,-1}+\delta_{m_22}-\delta_{m_2,-2})
\delta_{m_2,m'_2}\nn
&&\left (
\begin{array}{cccc}
0 & 1 & 0 & 0 \\
1 & 0 & 0 & 0 \\
0 & 0 & 1 & 0 \\
0 & 0 & 0 & 1 
\end{array} \right )
\end{eqnarray}

The character is given by the trace and is 

\begin{equation}
\chi^{T^1}(\sigma_d)=\Trace{T^1(\sigma_d)}=2
\end{equation}

% ------------------------------------------------------------------------------

\subsubsection{The xz mirror plane  $\sigma'_d$}

For the xz mirror plane $\sigma'_d$ the transformation is given by

\begin{eqnarray}
(x,y,z) &\stackrel{\sigma'_d}{\rightarrow}& (x,-y,z) \nn
(J_x,J_y,J_z) &\stackrel{\sigma'_d}{\rightarrow}&  (-J_x,J_y,-J_z) \nn
J_\pm &\stackrel{\sigma'_d}{\rightarrow}& - J_\mp \nn
\end{eqnarray}

\begin{eqnarray}
´O_{2-2}&\stackrel{\sigma'_d}{\rightarrow}&-O_{2-2} \nn
O_{2-1}&\stackrel{\sigma'_d}{\rightarrow}&-O_{2-1} \nn
O_{20}&\stackrel{\sigma'_d}{\rightarrow}&O_{20} \nn
O_{21}&\stackrel{\sigma'_d}{\rightarrow}&O_{21}\nn
O_{22}&\stackrel{\sigma'_d}{\rightarrow}&O_{22}\nn
\end{eqnarray}

or in summary
\begin{equation}
O_{2m}\stackrel{\sigma'_d}{\rightarrow} (\delta_{m0}+\delta_{m1}-\delta_{m,-1}+\delta_{m2}-\delta_{m,-2}) O_{2,m} 
\end{equation}

and 

\begin{eqnarray}
T^1_{m_1m_2b,m'_1m'_2b'}(\sigma'_d)&=&(\delta_{m_10}+\delta_{m_11}-\delta_{m_1,-1}+\delta_{m_12}-\delta_{m_1,-2})
\delta_{m_1,m'_1}\nn
&&(\delta_{m_20}+\delta_{m_21}-\delta_{m_2,-1}+\delta_{m_22}-\delta_{m_2,-2})
\delta_{m_2,m'_2}\nn
&&(\delta_{b1}\delta_{b'1}+\delta_{b2}\delta_{b'2}+\delta_{b3}\delta_{b'4}+\delta_{b4}\delta_{b'3})
\end{eqnarray}

Writing the $bb'$-dependence in Matrix notation we get

\begin{eqnarray}\label{sdd}
T^1_{m_1m_2,m'_1m'_2}(\sigma'_d)&=&(\delta_{m_10}+\delta_{m_11}-\delta_{m_1,-1}+\delta_{m_12}-\delta_{m_1,-2})
\delta_{m_1,m'_1}\nn
&&(\delta_{m_20}+\delta_{m_21}-\delta_{m_2,-1}+\delta_{m_22}-\delta_{m_2,-2})
\delta_{m_2,m'_2}\nn
&&\left (
\begin{array}{cccc}
1 & 0 & 0 & 0 \\
0 & 1 & 0 & 0 \\
0 & 0 & 0 & 1 \\
0 & 0 & 1 & 0 
\end{array} \right )
\end{eqnarray}

The character is given by the trace and is 

\begin{equation}
\chi^{T^1}(\sigma'_d)=\Trace{T^1(\sigma'_d)}=2
\end{equation}


% PPPPPPPPPPPPPPPPPPPPPPPPPPPPPPPPPPPPPPPPPPPPPPPPPPPPPPPPPPPPPPPPPPPPPPPPPPPPPPPP


\subsubsection{The Projection Operator}
The projection operator $P^{A_1}$ is given by equation (4.51) in~\cite{elliott79-1}.

\begin{equation}
P^{A_1}=\frac{1}{g} \sum_G \chi^{A_1}(G) T^1(G)
\end{equation}

Here $g=8$ denotes the number of symmetry elements in the point group $\bar 4 m 2$ and
$\chi^{A_1}(G)=1$, because $A_1$ is the unit representation.
Inserting equations (\ref{E}),(\ref{C2}),(\ref{S4}),(\ref{bS4}),\dots we get

\begin{eqnarray}\label{P}
P^{A_1}_{m_1m_2,m'_1m'_2}&=&\delta_{m_1,m'_1}\delta_{m_2,m'_2}
\left (
\begin{array}{cccc}
1 & 0 & 0 & 0 \\
0 & 1 & 0 & 0 \\
0 & 0 & 1 & 0 \\
0 & 0 & 0 & 1 
\end{array} \right )\nn
&&+(-1)^{m_1+m_2}\delta_{m_1,m'_1}\delta_{m_2,m'_2}
\left (
\begin{array}{cccc}
0 & 1 & 0 & 0 \\
1 & 0 & 0 & 0 \\
0 & 0 & 0 & 1 \\
0 & 0 & 1 & 0 
\end{array} \right )\nn
&&+(\delta_{m_10}+\delta_{m_11}-\delta_{m_1,-1}-\delta_{m_12}-\delta_{m_1,-2})
\delta_{(-1)^{m_1}m_1,m'_1} \nn
&&
(\delta_{m_20}+\delta_{m_21}-\delta_{m_2,-1}-\delta_{m_22}-\delta_{m_2,-2})
\delta_{(-1)^{m_2}m_2,m'_2} \nn
&&\left (
\begin{array}{cccc}
0 & 0 & 0 & 1 \\
0 & 0 & 1 & 0 \\
1 & 0 & 0 & 0 \\
0 & 1 & 0 & 0 
\end{array} \right ) \nn
&&+(\delta_{m_10}-\delta_{m_11}+\delta_{m_1,-1}-\delta_{m_12}-\delta_{m_1,-2})
\delta_{(-1)^{m_1}m_1,m'_1} \nn
&&
(\delta_{m_20}-\delta_{m_21}+\delta_{m_2,-1}-\delta_{m_22}-\delta_{m_2,-2})
\delta_{(-1)^{m_2}m_2,m'_2} \nn
&&\left (
\begin{array}{cccc}
0 & 0 & 1 & 0 \\
0 & 0 & 0 & 1 \\
0 & 1 & 0 & 0 \\
1 & 0 & 0 & 0 
\end{array} \right )\nn
&&+(\delta_{m_10}-\delta_{m_11}-\delta_{m_1,-1}-\delta_{m_12}+\delta_{m_1,-2})
\delta_{(-1)^{m_1}m_1,m'_1}\nn
&&(\delta_{m_20}-\delta_{m_21}-\delta_{m_2,-1}-\delta_{m_22}+\delta_{m_2,-2})
\delta_{(-1)^{m_2}m_2,m'_2}\nn
&&\left (
\begin{array}{cccc}
0 & 0 & 1 & 0 \\
0 & 0 & 0 & 1 \\
1 & 0 & 0 & 0 \\
0 & 1 & 0 & 0 
\end{array} \right )\nn
&&+(\delta_{m_10}+\delta_{m_11}+\delta_{m_1,-1}-\delta_{m_12}+\delta_{m_1,-2})
\delta_{(-1)^{m_1}m_1,m'_1}\nn
&&(\delta_{m_20}+\delta_{m_21}+\delta_{m_2,-1}-\delta_{m_22}+\delta_{m_2,-2})
\delta_{(-1)^{m_2}m_2,m'_2}\nn
&&\left (
\begin{array}{cccc}
0 & 0 & 0 & 1 \\
0 & 0 & 1 & 0 \\
0 & 1 & 0 & 0 \\
1 & 0 & 0 & 0 
\end{array} \right )\nn
&&+(\delta_{m_10}-\delta_{m_11}+\delta_{m_1,-1}+\delta_{m_12}-\delta_{m_1,-2})
\delta_{m_1,m'_1}\nn
&&(\delta_{m_20}-\delta_{m_21}+\delta_{m_2,-1}+\delta_{m_22}-\delta_{m_2,-2})
\delta_{m_2,m'_2}\nn
&&\left (
\begin{array}{cccc}
0 & 1 & 0 & 0 \\
1 & 0 & 0 & 0 \\
0 & 0 & 1 & 0 \\
0 & 0 & 0 & 1 
\end{array} \right )\nn
&&+(\delta_{m_10}+\delta_{m_11}-\delta_{m_1,-1}+\delta_{m_12}-\delta_{m_1,-2})
\delta_{m_1,m'_1}\nn
&&(\delta_{m_20}+\delta_{m_21}-\delta_{m_2,-1}+\delta_{m_22}-\delta_{m_2,-2})
\delta_{m_2,m'_2}\nn
&&\left (
\begin{array}{cccc}
1 & 0 & 0 & 0 \\
0 & 1 & 0 & 0 \\
0 & 0 & 0 & 1 \\
0 & 0 & 1 & 0 
\end{array} \right )
\end{eqnarray}



\subsubsection{The 13 independent terms in the NN and NNN quadrupolar Hamiltonian}\label{13ind}

The number of independent terms $m_{A_1}$ in the NN and NNN quadrupolar Hamiltonian can be
determined by the formula (4.28) in \cite{elliott79-1}:

\begin{equation}
m_{A_1}=\frac{1}{g}\sum_p c_p \chi^{(A_1)\star}_p \chi^{(T^1)}_p = \frac{1}{8} (100+2\times2)=13
\end{equation}

In order to determine the form of the interaction, we apply the projection operator
(\ref{P}) to some basis vectors $|mm'b\rangle$ ($=O_{2m}(\mbf J)O_{2m'}(\mbf J_b)$)
choosing various combinations of $m,m',b$ until all 13 terms have been found.

Not all terms are visible in the bulk properties of $TbPO_4$, because for some
terms the quadrupolar mean fields will sum up to zero. Nevertheless these so called 
''hidden'' terms may leave a trace in the dispersion of modes.
In the following we list all visible terms followed by the hidden terms in the
Hamiltonian:


\subparagraph{The $K^{\alpha}$ term}

Applying $P^{A_1}$ to $|001\rangle$ yields the $O_{20}$ interaction, which
corresponds to the $K^{\alpha}$ term in the model of Morin/Jens.

\begin{equation}
|001\rangle+|002\rangle+|003\rangle+|004\rangle =O_{20}(\mbf J)\sum_{b=1,2,3,4}O_{20}(\mbf J_b)
\end{equation}

\subparagraph{The $K^{\epsilon}$ or $O_{2\pm1}O_{2\pm1}$ term}

Applying $P^{A_1}$ to $|111\rangle$ yields 

\begin{equation}
|111\rangle+|112\rangle+|-1-13\rangle+|-1-14\rangle =
O_{21}(\mbf J)[O_{21}(\mbf J_1)+O_{21}(\mbf J_2)]+O_{2-1}(\mbf J)[O_{2-1}(\mbf J_3)+O_{2-1}(\mbf J_4)]
\end{equation}

\subparagraph{The $K^{\epsilon'}$ or $O_{2\pm1}O_{2\pm1}$ term}

Applying $P^{A_1}$ to $|-1-11\rangle$ yields 

\begin{equation}
|-1-11\rangle+|-1-12\rangle+|113\rangle+|114\rangle =
O_{2-1}(\mbf J)[O_{2-1}(\mbf J_1)+O_{2-1}(\mbf J_2)]+O_{21}(\mbf J)[O_{21}(\mbf J_3)+O_{21}(\mbf J_4)]
\end{equation}


\subparagraph{The $K^{\gamma}$ or $O_{22}O_{22}$ term}

Applying $P^{A_1}$ to $|221\rangle$ yields 

\begin{equation}
|221\rangle+|222\rangle+|223\rangle+|224\rangle =
O_{22}(\mbf J)[O_{22}(\mbf J_1)+O_{22}(\mbf J_2)+O_{22}(\mbf J_3)+O_{22}(\mbf J_4)]
\end{equation}


\subparagraph{The $K^{\delta}$ or $O_{2-2}O_{2-2}$ term}

Applying $P^{A_1}$ to $|-2-21\rangle$ yields 

\begin{equation}
|-2-21\rangle+|-2-22\rangle+|-2-23\rangle+|-2-24\rangle =
O_{2-2}(\mbf J)[O_{2-2}(\mbf J_1)+O_{2-2}(\mbf J_2)+O_{2-2}(\mbf J_3)+O_{2-2}(\mbf J_4)]
\end{equation}

\subparagraph{The hidden $O_{22}O_{2\pm1}$ term}

Applying $P^{A_1}$ to $|211\rangle$ yields 

\begin{equation}\label{nn211}
|211\rangle-|212\rangle+|2-13\rangle-|2-14\rangle =
O_{22}(\mbf J)[O_{21}(\mbf J_1)-O_{21}(\mbf J_2)+O_{2-1}(\mbf J_3)-O_{2-1}(\mbf J_4)]
\end{equation}
   

\subparagraph{The hidden $O_{2\pm1}O_{22}$ term}

Applying $P^{A_1}$ to $|121\rangle$ yields 

\begin{equation}\label{nn121}
|121\rangle-|122\rangle+|-123\rangle-|-124\rangle =
O_{21}(\mbf J)[O_{22}(\mbf J_1)-O_{22}(\mbf J_2)]+O_{2-1}(\mbf J)[O_{22}(\mbf J_3)-O_{22}(\mbf J_4)]
\end{equation}
   
\subparagraph{The hidden $O_{2-2}O_{2\pm1}$ term}

Applying $P^{A_1}$ to $|-2-11\rangle$ yields 

\begin{equation}
|-2-11\rangle-|-2-12\rangle-|-213\rangle+|-214\rangle =
O_{2-2}(\mbf J)[O_{2-1}(\mbf J_1)-O_{2-1}(\mbf J_2)-O_{21}(\mbf J_3)+O_{21}(\mbf J_4)]
\end{equation}
 

\subparagraph{The hidden $O_{2\pm1}O_{2-2}$ term}

Applying $P^{A_1}$ to $|-1-21\rangle$ yields 

\begin{equation}
|-1-21\rangle-|-1-22\rangle-|1-23\rangle+|1-24\rangle =
O_{2-1}(\mbf J)[O_{2-2}(\mbf J_1)-O_{2-2}(\mbf J_2)]-O_{21}(\mbf J)[O_{2-2}(\mbf J_3)-O_{2-2}(\mbf J_4)]
\end{equation}
   


  \subparagraph{The hidden $O_{20}O_{22}$ term}

Applying $P^{A_1}$ to $|021\rangle$ yields 

\begin{equation}
|021\rangle+|022\rangle-|023\rangle-|024\rangle =
O_{20}(\mbf J)[O_{22}(\mbf J_1)+O_{22}(\mbf J_2)-O_{22}(\mbf J_3)-O_{22}(\mbf J_4)]
\end{equation}
   

\subparagraph{The hidden $O_{20}O_{2\pm1}$ term}

Applying $P^{A_1}$ to $|011\rangle$ yields 

\begin{equation}
|011\rangle-|012\rangle-|0-13\rangle+|0-14\rangle =
O_{20}(\mbf J)[O_{21}(\mbf J_1)-O_{21}(\mbf J_2)-O_{2-1}(\mbf J_3)+O_{2-1}(\mbf J_4)]
\end{equation}
    

\subparagraph{The hidden $O_{22}O_{20}$ term}

Applying $P^{A_1}$ to $|201\rangle$ yields 

\begin{equation}
|201\rangle+|202\rangle-|203\rangle-|204\rangle =
O_{22}(\mbf J)[O_{20}(\mbf J_1)+O_{20}(\mbf J_2)-O_{20}(\mbf J_3)-O_{20}(\mbf J_4)]
\end{equation}


\subparagraph{The hidden $O_{2\pm1}O_{20}$ term}

Applying $P^{A_1}$ to $|101\rangle$ yields 

\begin{equation}
|101\rangle-|102\rangle -|-103\rangle+|-104\rangle =
O_{21}(\mbf J)[O_{20}(\mbf J_1)-O_{20}(\mbf J_2)]-
O_{2-1}(\mbf J)[O_{20}(\mbf J_3)-O_{20}(\mbf J_4)]
\end{equation}

\subsection{The reducible representation $T^3(G)$ for third nearest Tb-Tb neighbours}

The third nearest neighbour interactions transform according to the
same reducible representation $T^3(G)$. Similar to $T^1(G)$
we will in the following sections determine the representation matrices 
$T^3_{m_1m_2b,m'_1m'_2b'}(G)$ which describe the transformation
of the state $|m_1m_2b\rangle$ when applying a point symmetry 
operation $G$:

\begin{equation}
|m_1m_2b\rangle \stackrel{G}{\rightarrow} \sum_{m'_1m'_2b'} T^3_{m_1m_2b,m'_1m'_2b'}(G)|m'_1m'_2b'\rangle
\end{equation}
% ------------------------------------------------------------------------------

\subsubsection{The group element $E$}
 
For the identity $E$ the representation 

\begin{equation}
T^3_{m_1m_2b,m'_1m'_2b'}(E)=\delta_{m_1,m'_1}\delta_{m_2,m'_2}\delta_{bb'}
\end{equation}

Writing the $bb'$-dependence in Matrix notation we get

\begin{equation}\label{E3}
T^1_{m_1m_2,m'_1m'_2}(E)=\delta_{m_1,m'_1}\delta_{m_2,m'_2}
\left (
\begin{array}{ccccccccc}
1 & 0 & 0 & 0 & 0 & 0 & 0 & 0\\
0 & 1 & 0 & 0 & 0 & 0 & 0 & 0\\
0 & 0 & 1 & 0 & 0 & 0 & 0 & 0\\
0 & 0 & 0 & 1 & 0 & 0 & 0 & 0\\ 
0 & 0 & 0 & 0 & 1 & 0 & 0 & 0\\ 
0 & 0 & 0 & 0 & 0 & 1 & 0 & 0\\ 
0 & 0 & 0 & 0 & 0 & 0 & 1 & 0\\ 
0 & 0 & 0 & 0 & 0 & 0 & 0 & 1\\ 
\end{array} \right )
\end{equation}


The character is given by

\begin{equation}
\chi^{T^1}(E)=\Trace{T^1(E)}=5\times5\times8=200
\end{equation}

% ------------------------------------------------------------------------------

\subsubsection{The group element $C_2$}

We get

\begin{equation}\label{C23}
T^3_{m_1m_2,m'_1m'_2}(C_2)=(-1)^{m_1+m_2}\delta_{m_1,m'_1}\delta_{m_2,m'_2}
\left (
\begin{array}{ccccccccc}
0 & 0 & 0 & 0 & 0 & 0 & 1 & 0\\
0 & 0 & 0 & 0 & 0 & 0 & 0 & 1\\
0 & 0 & 0 & 0 & 1 & 0 & 0 & 0\\
0 & 0 & 0 & 0 & 0 & 1 & 0 & 0\\ 
0 & 0 & 1 & 0 & 0 & 0 & 0 & 0\\ 
0 & 0 & 0 & 1 & 0 & 0 & 0 & 0\\ 
1 & 0 & 0 & 0 & 0 & 0 & 0 & 0\\ 
0 & 1 & 0 & 0 & 0 & 0 & 0 & 0\\ 
\end{array} \right )
\end{equation}

The character is given by the trace and is zero

\begin{equation}
\chi^{T^3}(C_2)=\Trace{T^3(C_2)}=0
\end{equation}
% ------------------------------------------------------------------------------

\subsubsection{The clockwise mirror rotation $S_4$}

% (y,-x,-z)
 We get

\begin{eqnarray}\label{S43}
T^3_{m_1m_2,m'_1m'_2}(S_4)&=&(\delta_{m_10}-\delta_{m_11}+\delta_{m_1,-1}-\delta_{m_12}-\delta_{m_1,-2})
\delta_{(-1)^{m_1}m_1,m'_1} \nn
&&
(\delta_{m_20}-\delta_{m_21}+\delta_{m_2,-1}-\delta_{m_22}-\delta_{m_2,-2})
\delta_{(-1)^{m_2}m_2,m'_2} \nn
&&\left (
\begin{array}{ccccccccc}
0 & 0 & 0 & 1 & 0 & 0 & 0 & 0\\
0 & 0 & 1 & 0 & 0 & 0 & 0 & 0\\
0 & 0 & 0 & 0 & 0 & 0 & 0 & 1\\
0 & 0 & 0 & 0 & 0 & 0 & 1 & 0\\ 
0 & 1 & 0 & 0 & 0 & 0 & 0 & 0\\ 
1 & 0 & 0 & 0 & 0 & 0 & 0 & 0\\ 
0 & 0 & 0 & 0 & 0 & 1 & 0 & 0\\ 
0 & 0 & 0 & 0 & 1 & 0 & 0 & 0\\ 
\end{array} \right )
\end{eqnarray}

The character is given by the trace and is zero

\begin{equation}
\chi^{T^3}(S_4)=\Trace{T^1(S_4)}=0
\end{equation}

% ------------------------------------------------------------------------------

\subsubsection{The counter-clockwise mirror rotation $\bar S_4$}

% (-y,x,-z) \nn

We get

\begin{eqnarray}\label{bS43}
T^3_{m_1m_2,m'_1m'_2}(\bar S_4)&=&
(\delta_{m_10}+\delta_{m_11}-\delta_{m_1,-1}-\delta_{m_12}-\delta_{m_1,-2})
\delta_{(-1)^{m_1}m_1,m'_1} \nn
&&
(\delta_{m_20}+\delta_{m_21}-\delta_{m_2,-1}-\delta_{m_22}-\delta_{m_2,-2})
\delta_{(-1)^{m_2}m_2,m'_2} \nn
&&\left (
\begin{array}{ccccccccc}
0 & 0 & 0 & 0 & 0 & 1 & 0 & 0\\
0 & 0 & 0 & 0 & 1 & 0 & 0 & 0\\
0 & 1 & 0 & 0 & 0 & 0 & 0 & 0\\
1 & 0 & 0 & 0 & 0 & 0 & 0 & 0\\ 
0 & 0 & 0 & 0 & 0 & 0 & 0 & 1\\ 
0 & 0 & 0 & 0 & 0 & 0 & 1 & 0\\ 
0 & 0 & 0 & 1 & 0 & 0 & 0 & 0\\ 
0 & 0 & 1 & 0 & 0 & 0 & 0 & 0\\ 
\end{array} \right )
\end{eqnarray}

The character is given by the trace and is zero

\begin{equation}
\chi^{T^1}(\bar S_4)=\Trace{T^1(\bar S_4)}=0
\end{equation}

% ------------------------------------------------------------------------------

\subsubsection{The two fold horizontal [110]-axis  $C'_2$}

%(y,x,-z) 

We get

\begin{eqnarray}\label{bCd23}
T^3_{m_1m_2,m'_1m'_2}(C'_2)&=&(\delta_{m_10}-\delta_{m_11}-\delta_{m_1,-1}-\delta_{m_12}+\delta_{m_1,-2})
\delta_{(-1)^{m_1}m_1,m'_1}\nn
&&(\delta_{m_20}-\delta_{m_21}-\delta_{m_2,-1}-\delta_{m_22}+\delta_{m_2,-2})
\delta_{(-1)^{m_2}m_2,m'_2}\nn
&&\left (
\begin{array}{ccccccccc}
0 & 1 & 0 & 0 & 0 & 0 & 0 & 0\\
1 & 0 & 0 & 0 & 0 & 0 & 0 & 0\\
0 & 0 & 0 & 0 & 0 & 1 & 0 & 0\\
0 & 0 & 0 & 0 & 1 & 0 & 0 & 0\\ 
0 & 0 & 0 & 1 & 0 & 0 & 0 & 0\\ 
0 & 0 & 1 & 0 & 0 & 0 & 0 & 0\\ 
0 & 0 & 0 & 0 & 0 & 0 & 0 & 1\\ 
0 & 0 & 0 & 0 & 0 & 0 & 1 & 0\\ 
\end{array} \right )
\end{eqnarray}

The character is given by the trace and is zero

\begin{equation}
\chi^{T^1}(C'_2)=\Trace{T^1(C'_2)}=0
\end{equation}

% ------------------------------------------------------------------------------

\subsubsection{The two fold horizontal $[1\bar 10]$-axis  $C''_2$}

% (-y,-x,-z) 
We get

\begin{eqnarray}\label{bCdd23}
T^3_{m_1m_2,m'_1m'_2}(C''_2)&=&(\delta_{m_10}+\delta_{m_11}+\delta_{m_1,-1}-\delta_{m_12}+\delta_{m_1,-2})
\delta_{(-1)^{m_1}m_1,m'_1}\nn
&&(\delta_{m_20}+\delta_{m_21}+\delta_{m_2,-1}-\delta_{m_22}+\delta_{m_2,-2})
\delta_{(-1)^{m_2}m_2,m'_2}\nn
&&\left (
\begin{array}{ccccccccc}
0 & 0 & 0 & 0 & 0 & 0 & 0 & 1\\
0 & 0 & 0 & 0 & 0 & 0 & 1 & 0\\
0 & 0 & 0 & 1 & 0 & 0 & 0 & 0\\
0 & 0 & 1 & 0 & 0 & 0 & 0 & 0\\ 
0 & 0 & 0 & 0 & 0 & 1 & 0 & 0\\ 
0 & 0 & 0 & 0 & 1 & 0 & 0 & 0\\ 
0 & 1 & 0 & 0 & 0 & 0 & 0 & 0\\ 
1 & 0 & 0 & 0 & 0 & 0 & 0 & 0\\ 
\end{array} \right )
\end{eqnarray}

The character is given by the trace and is zero

\begin{equation}
\chi^{T^3}(C''_2)=\Trace{T^3(C''_2)}=0
\end{equation}

% ------------------------------------------------------------------------------

\subsubsection{The yz mirror plane  $\sigma_d$}

% (-x,y,z)

We get

\begin{eqnarray}\label{sd3}
T^3_{m_1m_2,m'_1m'_2}(\sigma_d)&=&(\delta_{m_10}-\delta_{m_11}+\delta_{m_1,-1}+\delta_{m_12}-\delta_{m_1,-2})
\delta_{m_1,m'_1}\nn
&&(\delta_{m_20}-\delta_{m_21}+\delta_{m_2,-1}+\delta_{m_22}-\delta_{m_2,-2})
\delta_{m_2,m'_2}\nn
&&\left (
\begin{array}{ccccccccc}
0 & 0 & 0 & 0 & 1 & 0 & 0 & 0\\
0 & 0 & 0 & 0 & 0 & 1 & 0 & 0\\
0 & 0 & 0 & 0 & 0 & 0 & 1 & 0\\
0 & 0 & 0 & 0 & 0 & 0 & 0 & 1\\ 
1 & 0 & 0 & 0 & 0 & 0 & 0 & 0\\ 
0 & 1 & 0 & 0 & 0 & 0 & 0 & 0\\ 
0 & 0 & 1 & 0 & 0 & 0 & 0 & 0\\ 
0 & 0 & 0 & 1 & 0 & 0 & 0 & 0\\ 
\end{array} \right )
\end{eqnarray}

The character is given by the trace and is zero

\begin{equation}
\chi^{T^3}(\sigma_d)=\Trace{T^3(\sigma_d)}=0
\end{equation}

% ------------------------------------------------------------------------------

\subsubsection{The xz mirror plane  $\sigma'_d$}

%(x,-y,z) 
We get

\begin{eqnarray}\label{sdd3}
T^3_{m_1m_2,m'_1m'_2}(\sigma'_d)&=&(\delta_{m_10}+\delta_{m_11}-\delta_{m_1,-1}+\delta_{m_12}-\delta_{m_1,-2})
\delta_{m_1,m'_1}\nn
&&(\delta_{m_20}+\delta_{m_21}-\delta_{m_2,-1}+\delta_{m_22}-\delta_{m_2,-2})
\delta_{m_2,m'_2}\nn
&&\left (
\begin{array}{ccccccccc}
0 & 0 & 1 & 0 & 0 & 0 & 0 & 0\\
0 & 0 & 0 & 1 & 0 & 0 & 0 & 0\\
1 & 0 & 0 & 0 & 0 & 0 & 0 & 0\\
0 & 1 & 0 & 0 & 0 & 0 & 0 & 0\\ 
0 & 0 & 0 & 0 & 0 & 0 & 1 & 0\\ 
0 & 0 & 0 & 0 & 0 & 0 & 0 & 1\\ 
0 & 0 & 0 & 0 & 1 & 0 & 0 & 0\\ 
0 & 0 & 0 & 0 & 0 & 1 & 0 & 0\\ 
\end{array} \right )
\end{eqnarray}

The character is given by the trace and is zero

\begin{equation}
\chi^{T^3}(\sigma'_d)=\Trace{T^3(\sigma'_d)}=0
\end{equation}


% PPPPPPPPPPPPPPPPPPPPPPPPPPPPPPPPPPPPPPPPPPPPPPPPPPPPPPPPPPPPPPPPPPPPPPPPPPPPPPPP


\subsubsection{The Projection Operator}
The projection operator $P^{A_1}$ is given by equation (4.51) in~\cite{elliott79-1}.

\begin{equation}
P^{A_1}=\frac{1}{g} \sum_G \chi^{A_1}(G) T^3(G)
\end{equation}

Here $g=8$ denotes the number of symmetry elements in the point group $\bar 4 m 2$ and
$\chi^{A_1}(G)=1$, because $A_1$ is the unit representation.
Inserting equations (\ref{E3}),(\ref{C23}),(\ref{S43}),(\ref{bS43}),\dots we get

\begin{eqnarray}\label{P3}
P^{A_1}_{m_1m_2,m'_1m'_2}&=&\delta_{m_1,m'_1}\delta_{m_2,m'_2}
\left (
\begin{array}{ccccccccc}
1 & 0 & 0 & 0 & 0 & 0 & 0 & 0\\
\dots 
\end{array} \right )\nn
&&+(\delta_{m_10}-\delta_{m_11}-\delta_{m_1,-1}-\delta_{m_12}+\delta_{m_1,-2})
\delta_{(-1)^{m_1}m_1,m'_1}\nn
&&(\delta_{m_20}-\delta_{m_21}-\delta_{m_2,-1}-\delta_{m_22}+\delta_{m_2,-2})
\delta_{(-1)^{m_2}m_2,m'_2}\nn
&&\left (
\begin{array}{ccccccccc}
0 & 1 & 0 & 0 & 0 & 0 & 0 & 0\\
\dots
\end{array} \right )\nn
&&+(\delta_{m_10}+\delta_{m_11}-\delta_{m_1,-1}+\delta_{m_12}-\delta_{m_1,-2})
\delta_{m_1,m'_1}\nn
&&(\delta_{m_20}+\delta_{m_21}-\delta_{m_2,-1}+\delta_{m_22}-\delta_{m_2,-2})
\delta_{m_2,m'_2}\nn
&&\left (
\begin{array}{ccccccccc}
0 & 0 & 1 & 0 & 0 & 0 & 0 & 0\\
\dots 
\end{array} \right )\nn
&&+(\delta_{m_10}-\delta_{m_11}+\delta_{m_1,-1}-\delta_{m_12}-\delta_{m_1,-2})
\delta_{(-1)^{m_1}m_1,m'_1} \nn
&&
(\delta_{m_20}-\delta_{m_21}+\delta_{m_2,-1}-\delta_{m_22}-\delta_{m_2,-2})
\delta_{(-1)^{m_2}m_2,m'_2} \nn
&&\left (
\begin{array}{ccccccccc}
0 & 0 & 0 & 1 & 0 & 0 & 0 & 0\\
\dots 
\end{array} \right ) \nn
&&+(\delta_{m_10}-\delta_{m_11}+\delta_{m_1,-1}+\delta_{m_12}-\delta_{m_1,-2})
\delta_{m_1,m'_1}\nn
&&(\delta_{m_20}-\delta_{m_21}+\delta_{m_2,-1}+\delta_{m_22}-\delta_{m_2,-2})
\delta_{m_2,m'_2}\nn
&&\left (
\begin{array}{ccccccccc}
0 & 0 & 0 & 0 & 1 & 0 & 0 & 0\\
\dots 
\end{array} \right )\nn
&&+(\delta_{m_10}+\delta_{m_11}-\delta_{m_1,-1}-\delta_{m_12}-\delta_{m_1,-2})
\delta_{(-1)^{m_1}m_1,m'_1} \nn
&&
(\delta_{m_20}+\delta_{m_21}-\delta_{m_2,-1}-\delta_{m_22}-\delta_{m_2,-2})
\delta_{(-1)^{m_2}m_2,m'_2} \nn
&&\left (
\begin{array}{ccccccccc}
0 & 0 & 0 & 0 & 0 & 1 & 0 & 0\\
\dots 
\end{array} \right )\nn
&&+(-1)^{m_1+m_2}\delta_{m_1,m'_1}\delta_{m_2,m'_2}
\left (
\begin{array}{ccccccccc}
0 & 0 & 0 & 0 & 0 & 0 & 1 & 0\\
\dots 
\end{array} \right )\nn
&&+(\delta_{m_10}+\delta_{m_11}+\delta_{m_1,-1}-\delta_{m_12}+\delta_{m_1,-2})
\delta_{(-1)^{m_1}m_1,m'_1}\nn
&&(\delta_{m_20}+\delta_{m_21}+\delta_{m_2,-1}-\delta_{m_22}+\delta_{m_2,-2})
\delta_{(-1)^{m_2}m_2,m'_2}\nn
&&\left (
\begin{array}{ccccccccc}
0 & 0 & 0 & 0 & 0 & 0 & 0 & 1\\
\dots 
\end{array} \right )
\end{eqnarray}


\subsubsection{The 25 independent terms in the third neighbour quadrupolar Hamiltonian}\label{25ind}

The number of independent terms $m_{A_1}$ in the third neighbour quadrupolar Hamiltonian can be
determined by the formula (4.28) in \cite{elliott79-1}:

\begin{equation}
m_{A_1}=\frac{1}{g}\sum_p c_p \chi^{(A_1)\star}_p \chi^{(T^3)}_p = \frac{1}{8} 200=25
\end{equation}

\subparagraph{The $K^{\alpha}$ term}

Applying $P^{A_1}$ to $|001\rangle$ yields the $O_{20}$ interaction, which
corresponds to the $K^{\alpha}$ term in the model of Morin/Jens.

\begin{equation}
O_{20}(\mbf J)\sum_{b=1,\dots,8}O_{20}(\mbf J_b)
\end{equation}

\subparagraph{The $K^{\epsilon}$ or $O_{2\pm1}O_{2\pm1}$ term}

Applying $P^{A_1}$ to $|111\rangle$ yields 

\begin{equation}
|111\rangle+|-1-12\rangle+|113\rangle+|-1-14\rangle+
|115\rangle+|-1-16\rangle+|117\rangle+|-1-18\rangle=
O_{21}(\mbf J)\sum_{b=1,3,5,7} O_{21}(\mbf J_b)+O_{2-1}(\mbf J)\sum_{b=2,4,6,8}O_{2-1}(\mbf J_b)
\end{equation}

\subparagraph{The $K^{\epsilon'}$ or $O_{2\pm1}O_{2\pm1}$ term}

Applying $P^{A_1}$ to $|-1-11\rangle$ yields 

\begin{equation}
|-1-11\rangle+|112\rangle+|-1-13\rangle+|114\rangle +
|-1-15\rangle+|116\rangle+|-1-17\rangle+|118\rangle =
O_{2-1}(\mbf J)\sum_{b=1,3,5,7} O_{2-1}(\mbf J_b)+O_{21}(\mbf J)\sum_{b=2,4,6,8}O_{21}(\mbf J_b)
\end{equation}


\subparagraph{The $K^{\gamma}$ or $O_{22}O_{22}$ term}

Applying $P^{A_1}$ to $|221\rangle$ yields 

\begin{equation}
|221\rangle+|222\rangle+|223\rangle+|224\rangle + \dots =
O_{22}(\mbf J)\sum_{b=1,\dots,8}O_{22}(\mbf J_b)
\end{equation}


\subparagraph{The $K^{\delta}$ or $O_{2-2}O_{2-2}$ term}

Applying $P^{A_1}$ to $|-2-21\rangle$ yields 

\begin{equation}
|-2-21\rangle+|-2-22\rangle+|-2-23\rangle+|-2-24\rangle + \dots=
O_{2-2}(\mbf J)\sum_{b=1,\dots,8}O_{2-2}(\mbf J_b)
\end{equation}


% ----------- hidden terms without m=1,-1 ---------------------

 \subparagraph{The hidden $O_{20}O_{22}$ term}

Applying $P^{A_1}$ to $|021\rangle$ yields 

\begin{eqnarray}
|021\rangle-|022\rangle+|023\rangle-|024\rangle
+|025\rangle-|026\rangle+|027\rangle-|028\rangle  =&& \nn
O_{20}(\mbf J)\sum_{b=1,...,8} (-1)^{b+1}O_{22}(\mbf J_b)&&
\end{eqnarray}
   

\subparagraph{The hidden $O_{22}O_{20}$ term}

Applying $P^{A_1}$ to $|201\rangle$ yields 

\begin{eqnarray}
|201\rangle-|202\rangle+|203\rangle-|204\rangle 
|205\rangle-|206\rangle+|207\rangle-|208\rangle 
=&&\nn
O_{22}(\mbf J)\sum_{b=1,...,8} (-1)^{b+1}O_{20}(\mbf J_b)&&
\end{eqnarray}

 \subparagraph{The hidden $O_{20}O_{2-2}$ term}

Applying $P^{A_1}$ to $|0-21\rangle$ yields 

\begin{eqnarray}
|0-21\rangle+|0-22\rangle-|0-23\rangle-|0-24\rangle
-|0-25\rangle-|0-26\rangle+|0-27\rangle+|0-28\rangle  =&& \nn
O_{20}(\mbf J)[
 O_{2-2}(\mbf J_1)+O_{2-2}(\mbf J_2)-O_{2-2}(\mbf J_3)-O_{2-2}(\mbf J_4)
-O_{2-2}(\mbf J_5)-O_{2-2}(\mbf J_6)+O_{2-2}(\mbf J_7)+O_{2-2}(\mbf J_8)
]&&
\end{eqnarray}
 

 \subparagraph{The hidden $O_{2-2}O_{20}$ term}

Applying $P^{A_1}$ to $|-201\rangle$ yields 

\begin{eqnarray}
|-201\rangle+|-202\rangle-|-203\rangle-|-204\rangle
-|-205\rangle-|-206\rangle+|-207\rangle+|-208\rangle  =&& \nn
O_{2-2}(\mbf J)[
 O_{20}(\mbf J_1)+O_{20}(\mbf J_2)-O_{20}(\mbf J_3)-O_{20}(\mbf J_4)
-O_{20}(\mbf J_5)-O_{20}(\mbf J_6)+O_{20}(\mbf J_7)+O_{20}(\mbf J_8)
]&&
\end{eqnarray}


 \subparagraph{The hidden $O_{2-2}O_{22}$ term}

Applying $P^{A_1}$ to $|-221\rangle$ yields 

\begin{eqnarray}
|-221\rangle-|-222\rangle-|-223\rangle+|-224\rangle
-|-225\rangle+|-226\rangle+|-227\rangle-|-228\rangle  =&& \nn
O_{2-2}(\mbf J)[
 O_{22}(\mbf J_1)-O_{22}(\mbf J_2)-O_{22}(\mbf J_3)+O_{22}(\mbf J_4)
-O_{22}(\mbf J_5)+O_{22}(\mbf J_6)+O_{22}(\mbf J_7)-O_{22}(\mbf J_8)
]&&
\end{eqnarray}

 \subparagraph{The hidden $O_{22}O_{2-2}$ term}

Applying $P^{A_1}$ to $|2-21\rangle$ yields 

\begin{eqnarray}
|2-21\rangle-|2-22\rangle-|2-23\rangle+|2-24\rangle
-|2-25\rangle+|2-26\rangle+|2-27\rangle-|2-28\rangle  =&& \nn
O_{22}(\mbf J)[
 O_{2-2}(\mbf J_1)-O_{2-2}(\mbf J_2)-O_{2-2}(\mbf J_3)+O_{2-2}(\mbf J_4)
-O_{2-2}(\mbf J_5)+O_{2-2}(\mbf J_6)+O_{2-2}(\mbf J_7)-O_{2-2}(\mbf J_8)
]&&
\end{eqnarray}


% ----------- hidden terms with m=1,-1 ---------------------

\subparagraph{The hidden $O_{2\pm1}O_{2\mp1}$ term}

Applying $P^{A_1}$ to $|1-11\rangle$ yields 

\begin{eqnarray}
|1-111\rangle+|-112\rangle-|1-13\rangle-|-114\rangle
-|1-15\rangle-|-116\rangle+|1-17\rangle+|-118\rangle =&&\nn
O_{21}(\mbf J)
[O_{2-1}(\mbf J_1)-O_{2-1}(\mbf J_3)-O_{2-1}(\mbf J_5)+O_{2-1}(\mbf J_7)]
+O_{2-1}(\mbf J)
[O_{21}(\mbf J_2)-O_{21}(\mbf J_4)-O_{21}(\mbf J_6)+O_{21}(\mbf J_8)]&&
\end{eqnarray}

\subparagraph{The hidden $O_{2\mp1}O_{2\pm1}$ term}

Applying $P^{A_1}$ to $|-111\rangle$ yields 

\begin{eqnarray}
|-1111\rangle+|1-12\rangle-|-113\rangle-|1-14\rangle
-|-115\rangle-|1-16\rangle+|-117\rangle+|1-18\rangle =&&\nn
O_{2-1}(\mbf J)
[O_{21}(\mbf J_1)-O_{21}(\mbf J_3)-O_{21}(\mbf J_5)+O_{21}(\mbf J_7)]
+O_{21}(\mbf J)
[O_{2-1}(\mbf J_2)-O_{2-1}(\mbf J_4)-O_{2-1}(\mbf J_6)+O_{2-1}(\mbf J_8)]&&
\end{eqnarray}

\subparagraph{The hidden $O_{20}O_{2\pm1}$ term}

Applying $P^{A_1}$ to $|011\rangle$ yields 

\begin{eqnarray}
|011\rangle-|0-12\rangle+|013\rangle-|0-14\rangle
-|015\rangle+|0-16\rangle-|017\rangle+|0-18\rangle =&&\nn
O_{20}(\mbf J)[
 O_{21}(\mbf J_1)-O_{2-1}(\mbf J_2)+O_{21}(\mbf J_3)-O_{2-1}(\mbf J_4)
-O_{21}(\mbf J_5)+O_{2-1}(\mbf J_6)-O_{21}(\mbf J_7)+O_{2-1}(\mbf J_8)
]&&
\end{eqnarray}

\subparagraph{The hidden $O_{2\pm1}O_{20}$ term}

Applying $P^{A_1}$ to $|101\rangle$ yields 

\begin{eqnarray}
 |101\rangle-|-102\rangle +|103\rangle-|-104\rangle 
-|105\rangle+|-106\rangle -|107\rangle+|-108\rangle 
=&&\nn
O_{21}(\mbf J)[O_{20}(\mbf J_1)+O_{20}(\mbf J_3)-O_{20}(\mbf J_5)-O_{20}(\mbf J_7)]-
O_{2-1}(\mbf J)[-O_{20}(\mbf J_2)-O_{20}(\mbf J_4)+O_{20}(\mbf J_6)+O_{20}(\mbf J_8)]&&
\end{eqnarray}


\subparagraph{The hidden $O_{20}O_{2\mp1}$ term}

Applying $P^{A_1}$ to $|0-11\rangle$ yields 

\begin{eqnarray}
|0-11\rangle-|012\rangle-|0-13\rangle+|014\rangle
+|0-15\rangle-|016\rangle-|0-17\rangle+|018\rangle =&&\nn
O_{20}(\mbf J)[
 O_{2-1}(\mbf J_1)-O_{21}(\mbf J_2)-O_{2-1}(\mbf J_3)+O_{21}(\mbf J_4)
+O_{2-1}(\mbf J_5)-O_{21}(\mbf J_6)-O_{2-1}(\mbf J_7)+O_{21}(\mbf J_8)
]&&
\end{eqnarray}

\subparagraph{The hidden $O_{2\mp1}O_{20}$ term}

Applying $P^{A_1}$ to $|-101\rangle$ yields 

\begin{eqnarray}
 |-101\rangle-|102\rangle -|-103\rangle+|104\rangle 
+|-105\rangle-|106\rangle -|-107\rangle+|108\rangle 
=&&\nn
O_{2-1}(\mbf J)[O_{20}(\mbf J_1)-O_{20}(\mbf J_3)+O_{20}(\mbf J_5)-O_{20}(\mbf J_7)]-
O_{21}(\mbf J) [O_{20}(\mbf J_2)-O_{20}(\mbf J_4)+O_{20}(\mbf J_6)-O_{20}(\mbf J_8)]&&
\end{eqnarray}



\subparagraph{The hidden $O_{22}O_{2\pm1}$ term}

Applying $P^{A_1}$ to $|211\rangle$ yields 

\begin{eqnarray}\label{nnn211}
|211\rangle+|2-12\rangle+|213\rangle+|2-14\rangle-
|215\rangle-|2-16\rangle-|217\rangle-|2-18\rangle&=& \nn
O_{22}(\mbf J)[O_{21}(\mbf J_1)+O_{2-1}(\mbf J_2)+O_{21}(\mbf J_3)+O_{2-1}(\mbf J_4)-
O_{21}(\mbf J_5)-O_{2-1}(\mbf J_6)-O_{21}(\mbf J_7)-O_{2-1}(\mbf J_8)]&&
\end{eqnarray}
   

\subparagraph{The hidden $O_{2\pm1}O_{22}$ term}

Applying $P^{A_1}$ to $|121\rangle$ yields 

\begin{eqnarray}\label{nnn121}
|121\rangle+|-122\rangle+|123\rangle+|-124\rangle 
-|125\rangle-|-126\rangle-|127\rangle-|-128\rangle=&& \nn
O_{21}(\mbf J)[O_{22}(\mbf J_1)+O_{22}(\mbf J_3)-O_{22}(\mbf J_5)-O_{22}(\mbf J_7)]+
O_{2-1}(\mbf J)[O_{22}(\mbf J_2)+O_{22}(\mbf J_4)-O_{22}(\mbf J_6)-O_{22}(\mbf J_8)]&&
\end{eqnarray}
   
\subparagraph{The hidden $O_{22}O_{2\mp1}$ term}

Applying $P^{A_1}$ to $|2-11\rangle$ yields 

\begin{eqnarray}
 |2-11\rangle+|212\rangle-|2-13\rangle-|214\rangle
+|2-15\rangle+|216\rangle-|2-17\rangle-|218\rangle&=& \nn
O_{22}(\mbf J)[O_{2-1}(\mbf J_1)+O_{21}(\mbf J_2)-O_{2-1}(\mbf J_3)-O_{21}(\mbf J_4)
+O_{2-1}(\mbf J_5)+O_{21}(\mbf J_6)-O_{2-1}(\mbf J_7)-O_{21}(\mbf J_8)]&&
\end{eqnarray}
   
\subparagraph{The hidden $O_{2\mp1}O_{22}$ term}

Applying $P^{A_1}$ to $|-121\rangle$ yields 

\begin{eqnarray}
|-121\rangle+|122\rangle-|-123\rangle-|124\rangle 
+|-125\rangle+|126\rangle-|-127\rangle-|128\rangle=&& \nn
O_{2-1}(\mbf J)[O_{22}(\mbf J_1)-O_{22}(\mbf J_3)+O_{22}(\mbf J_5)-O_{22}(\mbf J_7)]+
O_{21}(\mbf J)[O_{22}(\mbf J_2)-O_{22}(\mbf J_4)+O_{22}(\mbf J_6)-O_{22}(\mbf J_8)]&&
\end{eqnarray}

\subparagraph{The hidden $O_{2-2}O_{2\pm1}$ term}

Applying $P^{A_1}$ to $|-211\rangle$ yields 

\begin{eqnarray}
|-211\rangle-|-2-12\rangle-|-213\rangle+|-2-14\rangle 
+|-215\rangle-|-2-16\rangle-|-217\rangle+|-2-18\rangle =&& \nn
O_{2-2}(\mbf J)[O_{21}(\mbf J_1)-O_{2-1}(\mbf J_2)-O_{21}(\mbf J_3)+O_{2-1}(\mbf J_4)
+O_{21}(\mbf J_5)-O_{2-1}(\mbf J_6)-O_{21}(\mbf J_7)+O_{2-1}(\mbf J_8)]&&
\end{eqnarray}
 
\subparagraph{The hidden $O_{2\pm1}O_{2-2}$ term}

Applying $P^{A_1}$ to $|1-21\rangle$ yields 

\begin{eqnarray}
|1-21\rangle-|-1-22\rangle-|1-23\rangle+|-1-24\rangle 
+|1-25\rangle-|-1-26\rangle-|1-27\rangle+|-1-28\rangle =&& \nn
O_{21}(\mbf J)[O_{22}(\mbf J_1)-O_{22}(\mbf J_3)+O_{22}(\mbf J_5)-O_{22}(\mbf J_7)]
-O_{2-1}(\mbf J)[O_{22}(\mbf J_2)-O_{22}(\mbf J_4)+O_{22}(\mbf J_6)-O_{22}(\mbf J_8)]&&
\end{eqnarray}
 

\subparagraph{The hidden $O_{2-2}O_{2\mp1}$ term}

Applying $P^{A_1}$ to $|-2-11\rangle$ yields 

\begin{eqnarray}
|-2-11\rangle-|-212\rangle+|-2-13\rangle-|-214\rangle 
-|-2-15\rangle+|-216\rangle-|-2-17\rangle+|-218\rangle =&& \nn
O_{2-2}(\mbf J)[O_{2-1}(\mbf J_1)-O_{21}(\mbf J_2)+O_{2-1}(\mbf J_3)-O_{21}(\mbf J_4)
-O_{2-1}(\mbf J_5)+O_{21}(\mbf J_6)-O_{2-1}(\mbf J_7)+O_{21}(\mbf J_8)]&&
\end{eqnarray}
 

\subparagraph{The hidden $O_{2\mp1}O_{2-2}$ term}

Applying $P^{A_1}$ to $|-1-21\rangle$ yields 

\begin{eqnarray}
|-1-21\rangle-|1-22\rangle+|-1-23\rangle-|1-24\rangle 
-|-1-25\rangle+|1-26\rangle-|-1-27\rangle+|1-28\rangle =&&\nn
O_{2-1}(\mbf J)[O_{2-2}(\mbf J_1)+O_{2-2}(\mbf J_3)-O_{2-2}(\mbf J_5)-O_{2-2}(\mbf J_7)]
-O_{21}(\mbf J)[O_{2-2}(\mbf J_2)+O_{2-2}(\mbf J_4)-O_{2-2}(\mbf J_6)-O_{2-2}(\mbf J_8)]&&
\end{eqnarray}
   


